% Generated draft: requires mathematical and source review.
\documentclass[11pt]{article}
\usepackage[margin=25mm]{geometry}
\usepackage{fontspec}
\setmainfont{Latin Modern Roman}
\usepackage{amsmath,amssymb,mathtools,unicode-math}
\setmathfont{Latin Modern Math}
\usepackage{enumitem,needspace,xurl,hyperref}
\setlength{\parindent}{0pt}
\setlength{\parskip}{5pt}
\setlength{\emergencystretch}{3em}
\newcommand{\N}{\mathbb{N}}
\newcommand{\Z}{\mathbb{Z}}
\newcommand{\Q}{\mathbb{Q}}
\newcommand{\R}{\mathbb{R}}
\newcommand{\C}{\mathbb{C}}
\newcommand{\F}{\mathbb{F}}
\newcommand{\A}{\mathbb{A}}
\newcommand{\PP}{\mathbb{P}}
\newcommand{\E}{\mathbb{E}}
\newcommand{\eps}{\varepsilon}
\newcommand{\dd}{\,\mathrm{d}}
\begin{document}
\section*{UnsolvedMath ID: 20001906}
\textbf{Source problem code:} catalogue entry
\subsection{Definitions and mathematical statement}
The original Question 7.2.2 asks whether ``this version of the parallel postulate'' forces flatness, but omits its antecedent. The supplied, explicitly unverified reconstruction is the following Playfair-rigidity problem.

Let \((P,g)\) be a complete smooth Riemannian surface diffeomorphic to \(\R^2\). A complete geodesic image is the image of a nonconstant geodesic \(\gamma:\R\to P\), with parametrizations and orientation ignored. Assume
\[
\forall\ell\ \forall p\in P\setminus\ell,\quad
\#\{\text{complete geodesic images }m:p\in m,\ m\cap\ell=\varnothing\}=1.
\]
Must \(K_g\equiv0\), equivalently must \((P,g)\) be isometric to the Euclidean plane? No other synthetic Euclidean axioms, such as congruence axioms, are assumed.

The supplied research aid reports an affirmative answer with the additional no-focal-points hypothesis. One standard formulation of that hypothesis is that along every unit-speed geodesic, every nonzero normal Jacobi field \(J\) with \(J(0)=0\) satisfies \(\frac{\dd}{\dd t}|J(t)|^2>0\) for \(t>0\). The aid describes forward and backward asymptotic geodesics through an off-line point, attributed to O'Sullivan; Playfair uniqueness makes their images coincide, giving a bi-asymptotic pair, and a no-focal flat-strip theorem then supplies a Euclidean strip through the point. These are supplied proof ingredients, not independently checked theorems.

The aid separately reports flatness under \(K_g\geq0\), using Cheeger--Gromoll. It consequently describes any nonflat example satisfying the reconstructed axiom as having focal points and curvature of both signs; the nonpositive-curvature case is included in the no-focal obstruction. It also says that a counterexample must fail to admit total curvature, but neither the intended integrability convention nor a supporting reference is supplied. That last phrase is retained as an unresolved qualification, not converted into a particular assertion about \(\int_P|K|\,\dd A\). The unrestricted reconstructed question remains unresolved in the aid.
\subsection{Short English statement}
Does the unique-disjoint-geodesic version of Playfair's postulate force a complete Riemannian plane to be Euclidean? The supplied aid reports rigidity without focal points and under nonnegative curvature. The original extract omits the postulate itself, so the precise axiom is a supplied reconstruction.
\subsection{Sources}
\begin{itemize}
\item Question 7.2.2, [Bur-Kni1991], in the AIM geodesics problem list: \url{https://aimath.org/pastworkshops/geodesicsproblems.pdf}. The record supplies no full entry for that citation.
\item Local registry \texttt{data/unsolvedmath/problems.json}, permanent ID 20001906, code \texttt{catalogue entry}; canonical \texttt{aim-geometry-notes.json}, \texttt{notes[243]}, source item 7.2; updated 14 August 2026. Statement status reconstructed unverified; the original literature field is empty.
\item Supplied attempt 1: partial result, problem status at run open, rigor label proved, novelty confidence low; current registry status partially solved. O'Sullivan and Cheeger--Gromoll are named without bibliographic details. The aid is machine generated, not peer reviewed, and requires independent verification before citation.
\end{itemize}
\end{document}
